\documentclass{article}%
\usepackage[T1]{fontenc}%
\usepackage[utf8]{inputenc}%
\usepackage{lmodern}%
\usepackage{textcomp}%
\usepackage{lastpage}%
\usepackage{geometry}%
\geometry{margin=2.2cm}%
%
\usepackage{amsmath}%
\usepackage{amssymb}%
\usepackage{booktabs}%
\usepackage{longtable}%
\usepackage{array}%
\usepackage{hyperref}%
\hypersetup{hidelinks}%
\title{IME 2026 {-} Matemática (claude\_haiku\_5\_5) x Poliedro\\[4pt]\large Comparação com o gabarito Poliedro Resolve\thanks{Gerado com gpt-resolve: \url{https://github.com/lgabs/gpt-resolve}}}%
\date{10/11/2026}%
%
\begin{document}%
\normalsize%
\maketitle%
\section*{Resumo}\textbf{Nota média: 10,0/10} (10 de 10 itens com a resolução correta; cada item vale o mesmo dentro da questão).\par\medskip%
\begin{longtable}{@{}c c p{0.78\textwidth}@{}}\toprule Questão & Nota & Resumo \\ \midrule
1 & 10,0 & A resolução e o gabarito do Poliedro chegam à mesma resposta, $\alpha$ = 3 e $\beta$ = 2 (z = 3 + 2i), e ela está certa. Conferi direto: (3 + 2i)$^3$ = -9 + 46i, e a busca exaustiva nos inteiros mostra que a solução é única. A resolução usa o módulo ($\alpha$$^2$ + $\beta$$^2$ = 13), e o gabarito usa divisibilidade por 3 ($\alpha$ $\in$ \{$\pm$3, $\pm$9\}). Os dois argumentos estão completos e corretos. Não há divergência nem erro que mude a nota; só encontrei detalhes menores de redação. \\
2 & 10,0 & A resolução e o gabarito Poliedro dão a mesma resposta, x = -7/2 ou x = -23/2, e as duas estão corretas. Conferi com frações exatas que a equação equivale a (2x+7)$^3$(2x+23) = 0, o mesmo que (x + 7/2)$^{4}$ + (2x+7)$^3$ = 0. A raiz -7/2 é tripla e -23/2 é simples. O enunciado transcrito confere com o PDF original. Os dois desenvolvimentos são válidos, e não há divergência nem erro a apontar. \\
3 & 10,0 & A resolução está correta e bate com o gabarito do Poliedro: k = (1 - a)(1 + $\surd$a), que é o mesmo que (1 - $\surd$a)(1 + $\surd$a)$^2$. Conferi o enunciado no PDF original, refiz a conta e confirmei numericamente para vários valores de a, tanto o menor coeficiente angular 4/{[}(1 - $\surd$a)(1 + $\surd$a)$^2${]} quanto o ponto ótimo x* = (1 - $\surd$a)$^2$/4. A resolução é mais completa que o gabarito porque justifica a exclusão de a = 1/9 e o caso k = 0. O único deslize do gabarito é escrever 'k/4' no texto em vez de 4/k, sem efeito no resultado. \\
4 & 10,0 & A resolução está correta e equivale ao gabarito do Poliedro: as duas dão as mesmas três retas pela origem, sem a origem. A resolução escreve z = $\lambda$(t+i) com t $\in$ \{1, -2$\pm$$\surd$3\}. O gabarito escreve z = a(1+ri) com r $\in$ \{1, -2$\pm$$\surd$3\}. As formas coincidem porque (-2+$\surd$3)(-2-$\surd$3) = 1. Conferi o enunciado no PDF, refiz as fatorações à mão e testei as soluções numericamente em Python. Não há erro de conta em nenhum dos dois; só ficam observações de redação e de forma. \\
5 & 10,0 & Na questão 5, a resolução e o gabarito do Poliedro estão ambos corretos. O enunciado transcrito confere com o PDF original. A resolução prova a desigualdade escrevendo o produto menos 9 como (a-b)$^2$/(ab) + (a-c)$^2$/(ac) + (b-c)$^2$/(bc), soma de termos não negativos; conferi essa identidade com aritmética exata. Ela ainda mostra que a igualdade vale só quando a = b = c. O gabarito usa MA $\geq$ MG e prova essa desigualdade, para três termos, apenas numa observação opcional. Os dois caminhos são equivalentes e válidos; a prova da resolução é até mais completa, por ser autocontida. Não há erro de conta nem de lógica em nenhum dos dois. \\
6 & 10,0 & A resolução e o gabarito Poliedro chegam ao mesmo valor, 6 143 850, e ambos estão corretos. A resolução é até mais rigorosa que o gabarito: prova que a sequência tem de ser uma PA (a\_n = 3n - 2) e que a soma x\_j + y\_j + z\_j = a\_j - 3 vale para qualquer solução do sistema. Conferi as contas numericamente e não há divergência. \\
7 & 10,0 & A resolução e o gabarito do Poliedro dão a mesma resposta, P(I|N)=5/16=0,3125. Refiz a conta por Bayes e cheguei a 0,045/0,144=5/16, então os dois estão certos. O enunciado do PDF confere com a transcrição. A única falha que encontrei está no desenvolvimento do gabarito: o denominador traz 0,01$\cdot$0,99 no lugar de 0,1$\cdot$0,99. Isso não afeta a resposta final, que está correta. \\
8 & 10,0 & A resolução está correta e bate com o gabarito: a única solução real é (x, y, z) = (0, 0, 0). As contas que conferi estão certas: a derivada f' = 1 + 1/$\surd$(t$^2$+1) e os valores F($\pm$1) $\approx$ $\pm$4,17. O argumento do gabarito também vale, só está resumido e não prova que f é crescente. A resolução é mais rigorosa e completa: monotonicidade, equivalência com F(x) = 0, unicidade pelo Teorema do Valor Médio e um método alternativo igual ao do gabarito. O único ponto a levar a um professor é a base do 'log', que não está especificada, mas não muda a resposta. \\
9 & 10,0 & A resolução está certa e a resposta é equivalente à do gabarito Poliedro. As formas a($\surd$3-1)/(2+$\surd$2) e ($\surd$3-1)(2-$\surd$2)a/2 são iguais e valem $\approx$ 0,2144a. Conferi as contas numericamente e o enunciado no PDF original. A resolução ainda justifica que M e N são simétricos (AM = AN $\Rightarrow$ m = n), descarta a raiz t = (2+$\surd$3)a por ficar fora do lado e mostra que a configuração é única. Só há detalhes de notação, sem efeito na nota. \\
10 & 10,0 & A Questão 10 está correta e bate com o gabarito do Poliedro Resolve: V = $\pi$$\surd$6$\cdot$a$^3$/108 $\approx$ 0,0713$\cdot$a$^3$. As duas usam o mesmo raciocínio. Primeiro, r = a$\surd$6/12. Depois, o plano tangente passa pelos pontos médios das arestas e cria tetraedros de razão 1/2. A cada nível, o número de esferas vezes 4 e o volume de cada esfera vezes 1/8 dão razão 1/2, e a soma fica 2v0. Conferi as contas em Python e o enunciado com o PDF original. O único ponto de leitura é incluir ou não a esfera inscrita original. As duas a incluem, como o enunciado sugere, e a resolução avisa que, sem ela, o resultado seria a metade. \\
\bottomrule\end{longtable}%
\section*{Questão 1}%
\href{https://poliedroresolve.sistemapoliedro.com.br/vestibulares/ime/2026/ime-2-fase-dia-1-23-10-2025-2-fase-dia-1/dissertativa-1-matematica-2-fase-dia-1-ime-2026}{Resolução do gabarito}\par\smallskip%
\begin{longtable}{@{}p{0.06\textwidth} p{0.33\textwidth} p{0.33\textwidth} p{0.1\textwidth} p{0.08\textwidth}@{}}\toprule Item & Resolução & Gabarito & Veredito & Certo \\ \midrule
Q1 ($\alpha$ e $\beta$) & $\alpha$ = 3 e $\beta$ = 2 (z = 3 + 2i) & $\alpha$ = 3 e $\beta$ = 2 (z = 3 + 2i) & igual & ambos \\
\bottomrule\end{longtable}%
\begin{itemize}%
\item%
\textbf{Q1 ($\alpha$ e $\beta$))} Conferi o enunciado no PDF original (página 1) e ele é igual à transcrição. As duas respostas coincidem e estão certas. A busca exaustiva em python com $\alpha$, $\beta$ inteiros não nulos entre -20 e 20 encontra só (3, 2), e (3 + 2i)$^3$ = -9 + 46i. Os caminhos são diferentes, mas ambos são válidos. A resolução usa o módulo: ($\alpha$$^2$ + $\beta$$^2$)$^3$ = 2197 = 13$^3$, então $\alpha$$^2$ + $\beta$$^2$ = 13 e \{$\alpha$$^2$, $\beta$$^2$\} = \{4, 9\}. O caso $\alpha$$^2$ = 4 dá $\alpha$ = 9/23, que não é inteiro, e é eliminado. O caso $\alpha$$^2$ = 9 dá $\alpha$ = 3 e $\beta$ = 2. O gabarito mostra que 3 divide $\alpha$$^3$ e, portanto, divide $\alpha$. Assim, $\alpha$ $\in$ \{$\pm$3, $\pm$9\}, e só $\alpha$ = 3 dá $\beta$ inteiro ($\beta$ = $\pm$2). A equação (II) fixa $\beta$ = 2. Refiz os quatro casos do gabarito e estão certos: $\beta$$^2$ = 4, 2, 82/3 e 80/3. Também confirmei os sinais citados na resolução: (3, -2) $\to$ -9 - 46i, (-3, 2) $\to$ 9 + 46i e (-3, -2) $\to$ 9 - 46i.%
\end{itemize}%
\paragraph{Pontos de dúvida}%
\begin{itemize}%
\item%
\textit{(solution)} No passo 4, a resolução diz que $\alpha$$^2$ e $\beta$$^2$ são quadrados $\leq$ 13. Na verdade são $\leq$ 12, porque o outro termo vale pelo menos 1. Isso não altera nada. O passo 6, que confere os sinais das outras três combinações, é redundante: no caso B, as equações lineares -3$\alpha$ = -9 e 23$\beta$ = 46 já fixam $\alpha$ = 3 e $\beta$ = 2. Isso é só redação e não é erro.%
\item%
\textit{(key)} O gabarito escreve $\alpha$, $\beta$ $\in$ $\mathbb{Z}$ sem dizer "não nulos". Como nenhum candidato é nulo, isso não afeta o resultado. O passo "3 divide $\alpha$$^3$, então 3 divide $\alpha$" usa que 3 é primo sem dizer isso. É uma omissão aceitável.%
\end{itemize}%
\section*{Questão 2}%
\href{https://poliedroresolve.sistemapoliedro.com.br/vestibulares/ime/2026/ime-2-fase-dia-1-23-10-2025-2-fase-dia-1/dissertativa-2-matematica-2-fase-dia-1-ime-2026}{Resolução do gabarito}\par\smallskip%
\begin{longtable}{@{}p{0.06\textwidth} p{0.33\textwidth} p{0.33\textwidth} p{0.1\textwidth} p{0.08\textwidth}@{}}\toprule Item & Resolução & Gabarito & Veredito & Certo \\ \midrule
Q2 (soluções reais da equação) & S = \{-23/2, -7/2\}, ou seja, x = -11,5 ou x = -3,5 (-7/2 é raiz tripla de P, -23/2 é raiz simples) & x = -7/2 ou x = -23/2 & igual & ambos \\
\bottomrule\end{longtable}%
\begin{itemize}%
\item%
\textbf{Q2 (soluções reais da equação))} O enunciado transcrito confere com a página 1 do PDF. Refiz as contas com frações exatas. P(-7/2) = 0 e P(-23/2) = 0. A identidade P(x) = (2x+7)$^3$(2x+23)/16 vale em 40 pontos inteiros, o que basta, porque os dois lados são polinômios de grau 4. O agrupamento do gabarito também confere: x$^{4}$ + 14x$^3$ + (147/2)x$^2$ + (343/2)x + 7$^{4}$/16 = (x + 7/2)$^{4}$. O gabarito chega a isso pelo binômio de Newton. A resolução faz o mesmo em dois passos, com w = x$^2$ + 7x e s = 2x + 7, e chega a s$^{4}$/16 + s$^3$ = 0, que é a mesma equação do gabarito, (2x+7)$^{4}$/2$^{4}$ + (2x+7)$^3$ = 0. Também conferi as verificações numéricas da resolução. Em x = -23/2, os numeradores 279841, -65536, -340676, 155526 e -31556 estão corretos e somam -2401. Em x = -7/2, a substituição direta e as relações de Girard (soma -22, produto 7889/16) também conferem. As duas respostas são idênticas e estão corretas.%
\end{itemize}%
\paragraph{Pontos de dúvida}%
\begin{itemize}%
\item%
\textit{(solution)} A resolução diz que -7/2 é raiz tripla. Está certo, mas o enunciado só pede as soluções reais e o conjunto solução tem dois elementos, que é o que a resolução e o gabarito apresentam. O caminho do gabarito, reconhecer (x + 7/2)$^{4}$ pelo binômio de Newton, é mais curto que a dupla substituição w $\to$ s da resolução. Isso não afeta a nota.%
\end{itemize}%
\section*{Questão 3}%
\href{https://poliedroresolve.sistemapoliedro.com.br/vestibulares/ime/2026/ime-2-fase-dia-1-23-10-2025-2-fase-dia-1/dissertativa-3-matematica-2-fase-dia-1-ime-2026}{Resolução do gabarito}\par\smallskip%
\begin{longtable}{@{}p{0.06\textwidth} p{0.33\textwidth} p{0.33\textwidth} p{0.1\textwidth} p{0.08\textwidth}@{}}\toprule Item & Resolução & Gabarito & Veredito & Certo \\ \midrule
Q3: valor de k em função de a & k = (1 - a)(1 + $\surd$a) = (1 - $\surd$a)(1 + $\surd$a)$^2$; secante mínima liga (a, f(a)) a x* = (1 - $\surd$a)$^2$/4, com inclinação 4/{[}(1 - $\surd$a)(1 + $\surd$a)$^2${]} & k = (1 - a)(1 + $\surd$a); menor coeficiente angular 4/{[}(1 - $\surd$a)(1 + $\surd$a)$^2${]}, igualado a 4/k & igual & ambos \\
\bottomrule\end{longtable}%
\begin{itemize}%
\item%
\textbf{Q3: valor de k em função de a)} O enunciado transcrito confere com o PDF original. Refiz a conta: com t = $\surd$x e s = $\surd$a, f(x) = t/(1 - t), e o coeficiente angular da secante é m = 1/{[}(1 - s)(1 - t)(t + s){]}. O produto (1 - t)(t + s) é máximo em t* = (1 - s)/2 e vale (1 + s)$^2$/4. Logo m\_min = 4/{[}(1 - s)(1 + s)$^2${]}. Como a reta ky - 4x = 2026 tem inclinação 4/k, sai k = (1 - s)(1 + s)$^2$ = (1 - a)(1 + $\surd$a). Fiz uma varredura numérica em x para a $\in$ \{0; 0,04; 0,25; 0,36; 0,5; 0,81; 0,95\}: o mínimo das inclinações coincide com 4/k e o ponto de mínimo coincide com x* = (1 - $\surd$a)$^2$/4. A resolução também mostra por que a = 1/9 é excluído: nesse caso t* = s, e o ínfimo vira a inclinação da tangente, que não é atingido por nenhuma secante. As resoluções diferem só no método, já que o gabarito usa o vértice/$\Delta$ de uma quadrática em $\surd$x e a resolução completa o quadrado. A resposta final é a mesma.%
\end{itemize}%
\paragraph{Pontos de dúvida}%
\begin{itemize}%
\item%
\textit{(key)} O texto do gabarito do Poliedro diz que o coeficiente angular de ky - 4x = 2026 é 'k/4'. O correto é 4/k, e é 4/k que a própria conta usa na linha seguinte. É erro de digitação e não muda a resposta.%
\item%
\textit{(key)} A resolução do gabarito escreve o ponto como '($\alpha$, \ldots{})' em vez de '(a, \ldots{})', e a fração da inclinação aparece com os parênteses mal transcritos. São só problemas de redação ou digitação.%
\item%
\textit{(key)} O gabarito não comenta por que a $\neq$ 1/9 nem verifica que o vértice t* = (1 - $\surd$a)/2 é diferente de $\surd$a, isto é, que o ponto ótimo é de fato uma secante. A resolução trata isso explicitamente. Fica como observação de rigor, que pode pesar na correção de uma prova discursiva.%
\end{itemize}%
\section*{Questão 4}%
\href{https://poliedroresolve.sistemapoliedro.com.br/vestibulares/ime/2026/ime-2-fase-dia-1-23-10-2025-2-fase-dia-1/dissertativa-4-matematica-2-fase-dia-1-ime-2026}{Resolução do gabarito}\par\smallskip%
\begin{longtable}{@{}p{0.06\textwidth} p{0.33\textwidth} p{0.33\textwidth} p{0.1\textwidth} p{0.08\textwidth}@{}}\toprule Item & Resolução & Gabarito & Veredito & Certo \\ \midrule
Q4 (valores de z) & z = $\lambda$(1+i), z = $\lambda$($\surd$3-2+i) ou z = $\lambda$(-2-$\surd$3+i), $\lambda$ $\in$ $\mathbb{R}$*. Equivalentemente, arg z $\in$ \{$\pi$/4, 7$\pi$/12, 11$\pi$/12\} (mod $\pi$): três retas pela origem, sem a origem. & z = a(1+i) ou z = a{[}1+(-2$\pm$$\surd$3)i{]}, a $\in$ $\mathbb{R}$* & equivalente & ambos \\
\bottomrule\end{longtable}%
\begin{itemize}%
\item%
\textbf{Q4 (valores de z))} Conferi o enunciado no PDF e ele bate com a transcrição. As duas resoluções chegam à mesma cúbica homogênea. Com z = x+yi, a resolução usa t = x/y e obtém t$^3$+3t$^2$-3t-1 = 0. O gabarito multiplica por 2ab$^2$ e obtém a$^3$+3a$^2$b-3ab$^2$-b$^3$ = 0, que é a mesma equação. Refiz à mão: (t+1)(t$^2$+2t-1) - 4t = t$^3$+3t$^2$-3t-1 = (t-1)(t$^2$+4t+1), e (a-b)(a$^2$+4ab+b$^2$) = a$^3$+3a$^2$b-3ab$^2$-b$^3$. A resolução fixa x/y $\in$ \{1, -2$\pm$$\surd$3\}; o gabarito fixa y/x $\in$ \{1, -2$\pm$$\surd$3\}. O conjunto \{1, -2+$\surd$3, -2-$\surd$3\} é fechado por inversão, porque (-2+$\surd$3)(-2-$\surd$3) = 1, então as retas são as mesmas. Os sinais aparecem trocados: a reta $\lambda$($\surd$3-2+i) da resolução é a reta a(1+(-2-$\surd$3)i) do gabarito, e vice{-}versa. Testei numericamente em Python os três tipos de z, para vários $\lambda$, nas duas formas: em todos o lado esquerdo vale 1. Para z = 2+3i, que não é solução, o valor dá $\approx$ -0,028. Os ângulos 45$^\circ$, 105$^\circ$ e 165$^\circ$ conferem: cot 105$^\circ$ = -2+$\surd$3 e cot 165$^\circ$ = -2-$\surd$3. As duas excluem a origem ($\lambda$ $\neq$ 0 e a $\neq$ 0) e os casos com xy = 0, em que os denominadores se anulam.%
\end{itemize}%
\paragraph{Pontos de dúvida}%
\begin{itemize}%
\item%
\textit{(key)} O gabarito não declara as condições de existência (a $\neq$ 0 e b $\neq$ 0, para que z-z? e z$^2$-z?$^2$ não se anulem). Elas ficam implícitas: o gabarito divide por a e por b e exige a $\in$ $\mathbb{R}$*. Também não faz falta excluir b = 0 à parte, porque com a $\neq$ 0 todas as soluções têm b = a$\cdot$r com r $\neq$ 0. A transcrição tem ainda uma falha de formatação ('= 1 a/b + ...', com duas linhas coladas), sem efeito no resultado.%
\item%
\textit{(solution)} A resolução escreve as soluções com x/y = -2$\pm$$\surd$3, e o gabarito com y/x = -2$\pm$$\surd$3. Os conjuntos são iguais, mas os sinais $\pm$ correspondem a retas trocadas: ($\surd$3-2+i) ? (1+(-2-$\surd$3)i). Quem comparar só a forma pode achar que as respostas diferem. Na forma polar, a resolução usa '$\rho$ $\in$ $\mathbb{R}$*' junto com '$\theta$ mod $\pi$', o que é redundante, mas está correto.%
\end{itemize}%
\section*{Questão 5}%
\href{https://poliedroresolve.sistemapoliedro.com.br/vestibulares/ime/2026/ime-2-fase-dia-1-23-10-2025-2-fase-dia-1/dissertativa-5-matematica-2-fase-dia-1-ime-2026}{Resolução do gabarito}\par\smallskip%
\begin{longtable}{@{}p{0.06\textwidth} p{0.33\textwidth} p{0.33\textwidth} p{0.1\textwidth} p{0.08\textwidth}@{}}\toprule Item & Resolução & Gabarito & Veredito & Certo \\ \midrule
Q5 -- demonstração de (a+b+c)(1/a+1/b+1/c) $\geq$ 9 para a, b, c > 0 & Prova direta. Expandindo o produto, obtém{-}se 3 + (a/b+b/a) + (a/c+c/a) + (b/c+c/b). Daí vem a identidade produto - 9 = (a-b)$^2$/(ab) + (a-c)$^2$/(ac) + (b-c)$^2$/(bc) $\geq$ 0, logo o produto é $\geq$ 9. A igualdade vale se e somente se a = b = c. Nas observações, a resolução também dá a prova por MA $\geq$ MG. & Prova por MA $\geq$ MG: a+b+c $\geq$ 3?(abc) e 1/a+1/b+1/c $\geq$ 3/?(abc). Como os dois lados são positivos, multiplicam{-}se as desigualdades e obtém{-}se $\geq$ 9 (c.q.d.). O gabarito também prova MA $\geq$ MG para três termos pela fatoração p$^3$+q$^3$+r$^3$-3pqr = (p+q+r)$\cdot$$\frac12${[}(p-q)$^2$+(p-r)$^2$+(q-r)$^2${]}. & equivalente & ambos \\
\bottomrule\end{longtable}%
\begin{itemize}%
\item%
\textbf{Q5 -- demonstração de (a+b+c)(1/a+1/b+1/c) $\geq$ 9 para a, b, c > 0)} As duas chegam à mesma conclusão por caminhos diferentes e ambas estão corretas. Conferi com aritmética racional exata em 10.000 casos aleatórios: a identidade da resolução, produto - 9 = $\Sigma$(x-y)$^2$/(xy), está certa. A fatoração do gabarito, p$^3$+q$^3$+r$^3$-3pqr = (p+q+r)$\cdot$$\frac12$$\Sigma$(p-q)$^2$, também. A prova da resolução é autocontida, porque não depende de MA $\geq$ MG sem demonstração, e ainda determina o caso de igualdade (a = b = c), que o enunciado não pede. Contas auxiliares conferidas: (1,2,3) dá 6$\cdot$11/6 = 11 = 9 + 1/2 + 4/3 + 1/6. a = b = c = 1 dá 9. O contraexemplo (-1,1,1) dá 1 < 9, o que mostra que a positividade é necessária.%
\end{itemize}%
\paragraph{Pontos de dúvida}%
\begin{itemize}%
\item%
\textit{(key)} A prova principal do gabarito usa MA $\geq$ MG para três termos sem demonstrá{-}la. A demonstração só aparece numa observação opcional ("caso o estudante tivesse tempo"). Se a banca exigir a prova completa, uma resposta que só cite MA $\geq$ MG pode perder pontos. A prova principal da resolução não tem esse problema, porque é elementar e autocontida.%
\item%
\textit{(solution)} A prova alternativa por MA $\geq$ MG, nas observações da resolução, também cita a desigualdade sem demonstrá{-}la. Isso não compromete nada, porque a prova principal (soma de quadrados sobre denominadores positivos) é completa.%
\end{itemize}%
\section*{Questão 6}%
\href{https://poliedroresolve.sistemapoliedro.com.br/vestibulares/ime/2026/ime-2-fase-dia-1-23-10-2025-2-fase-dia-1/dissertativa-6-matematica-2-fase-dia-1-ime-2026}{Resolução do gabarito}\par\smallskip%
\begin{longtable}{@{}p{0.06\textwidth} p{0.33\textwidth} p{0.33\textwidth} p{0.1\textwidth} p{0.08\textwidth}@{}}\toprule Item & Resolução & Gabarito & Veredito & Certo \\ \midrule
Questão 6 (soma de x\_j + y\_j + z\_j) & 6 143 850 (PA a\_n = 3n - 2; y\_j = 3; x\_j + y\_j + z\_j = 3j - 5 = a\_j - 3) & 6 143 850 (PA de razão 3; y\_j = 3; x\_j + y\_j + z\_j = a\_j - 3; soma = (1 + 6073)$\cdot$2025/2 - 6075) & igual & ambos \\
\bottomrule\end{longtable}%
\begin{itemize}%
\item%
\textbf{Questão 6 (soma de x\_j + y\_j + z\_j))} Conferi o enunciado com o PDF original (p. 2) e a transcrição está fiel. Refiz as contas em python3: $\Sigma$(3j - 5) para j = 1..2025 dá 6 143 850. Pelo caminho do gabarito, $\Sigma$a\_j - 3$\cdot$2025 = 6 149 925 - 6 075 = 6 143 850, o mesmo valor. A resolução prova que a condição de compatibilidade a\_j - 2a\_\{j+1\} + a\_\{j+2\} = 0 é necessária e também suficiente. Disso sai a PA de razão 3, com y\_j = 3 único e x\_j + z\_j = a\_j - 6, de modo que a soma não depende do parâmetro livre. As contas intermediárias também batem: 2025$\cdot$1013 = 2 051 325, vezes 3 dá 6 153 975, e 6 153 975 - 10 125 = 6 143 850. A conferência com j = 1 (x = -5, y = 3, z = 0) satisfaz as três equações.%
\end{itemize}%
\paragraph{Pontos de dúvida}%
\begin{itemize}%
\item%
\textit{(key)} O gabarito diz que 'uma solução possível' é a sequência ser PA. Essa redação sugere que a PA seria só uma das possibilidades. Na verdade, ela é obrigatória: subtraindo as equações, y = a\_\{j+1\} - a\_j = a\_\{j+2\} - a\_\{j+1\} para todo j, e isso força razão constante. A redação não muda o resultado, mas é imprecisa. A resolução conferida trata isso corretamente, mostrando que a condição é necessária e suficiente.%
\item%
\textit{(statement)} A matriz dos coeficientes tem posto 2, então o sistema tem infinitas soluções e (x\_j, y\_j, z\_j) não é única. A resposta só fica bem definida porque y\_j = 3 e x\_j + z\_j = a\_j - 6 são fixos. A resolução justifica isso explicitamente no Passo 4. O gabarito só faz isso de forma implícita.%
\end{itemize}%
\section*{Questão 7}%
\href{https://poliedroresolve.sistemapoliedro.com.br/vestibulares/ime/2026/ime-2-fase-dia-1-23-10-2025-2-fase-dia-1/dissertativa-7-matematica-2-fase-dia-1-ime-2026}{Resolução do gabarito}\par\smallskip%
\begin{longtable}{@{}p{0.06\textwidth} p{0.33\textwidth} p{0.33\textwidth} p{0.1\textwidth} p{0.08\textwidth}@{}}\toprule Item & Resolução & Gabarito & Veredito & Certo \\ \midrule
Q7 --- P(infectado | antivírus respondeu 'não infectado') & 5/16 = 0,3125 = 31,25\% & 5/16 & igual & ambos \\
\bottomrule\end{longtable}%
\begin{itemize}%
\item%
\textbf{Q7 --- P(infectado | antivírus respondeu 'não infectado'))} Refiz a conta por Bayes: P(N|I)=0,05, P(N|?)=0,99, P(I)=0,9 e P(?)=0,1. Daí P(N)=0,045+0,099=0,144 e P(I|N)=0,045/0,144=5/16. Conferi em python com frações exatas e o resultado é 5/16. O enunciado do PDF bate com a transcrição. As verificações extras da resolução também estão corretas: razão de chances 9$\cdot$5/99=5/11, contagem com 1000 computadores (45 de 144) e sensibilidade com p=0,5, que dá $\approx$4,8\%.%
\end{itemize}%
\paragraph{Pontos de dúvida}%
\begin{itemize}%
\item%
\textit{(key)} Há um erro de digitação no denominador da resolução do gabarito: está escrito '0,9$\cdot$0,05 + 0,01$\cdot$0,99', mas o correto é '0,9$\cdot$0,05 + 0,1$\cdot$0,99'. Com 0,01, a conta daria 0,045/0,0549 $\approx$ 0,82, não 5/16. A resposta final do gabarito (5/16) está certa; só a expressão intermediária está errada. Além disso, as barras de complemento se perderam na transcrição: 'P(B/A)' aparece três vezes, no lugar de P(B|A), P(B|?) e P(B?|A). Isso não muda a nota.%
\item%
\textit{(solution)} A resolução trata 'não indica vírus' como o complementar de 'indica vírus' (resposta binária, P(N|I)=1-0,95). É a leitura padrão e a mesma do gabarito, e a resolução declara essa hipótese. A seção de hipóteses e a análise de sensibilidade vão além do que a questão pede, mas não contêm erro.%
\end{itemize}%
\section*{Questão 8}%
\href{https://poliedroresolve.sistemapoliedro.com.br/vestibulares/ime/2026/ime-2-fase-dia-1-23-10-2025-2-fase-dia-1/dissertativa-8-matematica-2-fase-dia-1-ime-2026}{Resolução do gabarito}\par\smallskip%
\begin{longtable}{@{}p{0.06\textwidth} p{0.33\textwidth} p{0.33\textwidth} p{0.1\textwidth} p{0.08\textwidth}@{}}\toprule Item & Resolução & Gabarito & Veredito & Certo \\ \midrule
Q8 (provar que a solução é única e calculá{-}la) & Solução única (x, y, z) = (0, 0, 0). A resolução mostra que f(t) = t + ln(t + $\surd$(t$^2$+1)) tem f'(t) = 1 + 1/$\surd$(t$^2$+1) > 1. Daí F(t) = f(f(f(t))) - t tem F' > 0, logo no máximo uma raiz, e F(0) = 0. Também apresenta um método alternativo, sem derivadas, que conclui x = y = z = 0. & Solução única (x, y, z) = (0, 0, 0). O gabarito usa que f é estritamente crescente: de x = f(f(f(x))) segue f(x) = x, logo x = y = z. Então log(x + $\surd$(x$^2$+1)) = 0, que dá $\surd$(x$^2$+1) = 1 - x, e daí x = 0. & igual & ambos \\
\bottomrule\end{longtable}%
\begin{itemize}%
\item%
\textbf{Q8 (provar que a solução é única e calculá{-}la))} A conclusão é a mesma nas duas: a solução (0,0,0) é única. Conferi as contas. A derivada de ln(t + $\surd$(t$^2$+1)) é 1/$\surd$(t$^2$+1), então f' = 1 + 1/$\surd$(t$^2$+1) > 1 está certo. Numericamente F(1) $\approx$ 4,1714 e F(-1) $\approx$ -4,1714, como a resolução declara. Com log decimal, F($\pm$1) $\approx$ $\pm$1,474 e F(0) = 0. O argumento do gabarito vale, embora resumido: se f(x) > x, então f(f(x)) > f(x) > x e f(f(f(x))) > f(f(x)) > x, um absurdo; o caso f(x) < x é análogo. Ao elevar ao quadrado $\surd$(x$^2$+1) = 1 - x sai x = 0, que satisfaz 1 - x $\geq$ 0. A resolução é mais completa: prova que f é crescente, a equivalência do sistema com F(x) = 0, a unicidade pelo Teorema do Valor Médio e a independência da base do logaritmo.%
\end{itemize}%
\paragraph{Pontos de dúvida}%
\begin{itemize}%
\item%
\textit{(statement)} O enunciado escreve 'log' sem dizer a base. A resolução usa ln e mostra que o resultado vale para qualquer base b > 1, porque f\_b(0) = 0 e f\_b' > 1. O gabarito não trata disso. Isso não muda a resposta (0,0,0).%
\item%
\textit{(key)} O gabarito diz que f é estritamente crescente sem justificar. A justificativa é simples: t e ln(t + $\surd$(t$^2$+1)) = arcsenh(t) são crescentes. O gabarito também omite os passos intermediários de 'f(x) > x $\Rightarrow$ f(f(f(x))) > x', que exigem aplicar a monotonicidade duas vezes. Em uma prova de 'Prove que', a banca pode cobrar essas justificativas, mas a conclusão do gabarito está correta.%
\end{itemize}%
\section*{Questão 9}%
\href{https://poliedroresolve.sistemapoliedro.com.br/vestibulares/ime/2026/ime-2-fase-dia-1-23-10-2025-2-fase-dia-1/dissertativa-9-matematica-2-fase-dia-1-ime-2026}{Resolução do gabarito}\par\smallskip%
\begin{longtable}{@{}p{0.06\textwidth} p{0.33\textwidth} p{0.33\textwidth} p{0.1\textwidth} p{0.08\textwidth}@{}}\toprule Item & Resolução & Gabarito & Veredito & Certo \\ \midrule
Q9 --- raio do círculo inscrito no triângulo CMN & r = ($\surd$3-1)(2-$\surd$2)a/2 = (2$\surd$3+$\surd$2-$\surd$6-2)a/2 $\approx$ 0,214a & r = a($\surd$3 - 1)/(2 + $\surd$2) & equivalente & ambos \\
\bottomrule\end{longtable}%
\begin{itemize}%
\item%
\textbf{Q9 --- raio do círculo inscrito no triângulo CMN)} As duas formas são a mesma porque 1/(2+$\surd$2) = (2-$\surd$2)/2, já que (2+$\surd$2)(2-$\surd$2) = 2. Refiz em Python com a = 1. Com t = 2-$\surd$3 obtive AM = AN = MN $\approx$ 1,03528, o que confirma que o triângulo é equilátero. S/p calculado direto pelas coordenadas de C, M e N deu 0,2144127. As expressões da resolução, a forma expandida e a do gabarito dão o mesmo valor. O caminho também é igual nos dois: BM = a$\cdot$tg15$^\circ$ = a(2-$\surd$3) equivale a t = (2-$\surd$3)a; CMN é retângulo isósceles com catetos ($\surd$3-1)a; e r = S/p. Conferi o enunciado no PDF original (pág. 2) e ele bate com a transcrição.%
\end{itemize}%
\paragraph{Pontos de dúvida}%
\begin{itemize}%
\item%
\textit{(solution)} No item 5, a resolução escreve S = (r/2)(x+y+z) sem dizer que x, y, z são os lados. Na conferência ela também usa p para um cateto em r = (p+q-h)/2, sendo que p já era o semiperímetro. É só redação: não afeta o resultado.%
\item%
\textit{(key)} Na resolução do Poliedro, o quociente final aparece como (a$^2$($\surd$3-1)$^2$)/(2+$\surd$2)a($\surd$3-1), já sem os fatores 1/2 de S e de p, que se cancelam, e com os parênteses do denominador mal formatados. O resultado final está certo.%
\end{itemize}%
\section*{Questão 10}%
\href{https://poliedroresolve.sistemapoliedro.com.br/vestibulares/ime/2026/ime-2-fase-dia-1-23-10-2025-2-fase-dia-1/dissertativa-10-matematica-2-fase-dia-1-ime-2026}{Resolução do gabarito}\par\smallskip%
\begin{longtable}{@{}p{0.06\textwidth} p{0.33\textwidth} p{0.33\textwidth} p{0.1\textwidth} p{0.08\textwidth}@{}}\toprule Item & Resolução & Gabarito & Veredito & Certo \\ \midrule
Q10 --- volume total das esferas interiores ao tetraedro ABCD & V = ($\pi$$\surd$6/108)$\cdot$a$^3$ $\approx$ 0,0713$\cdot$a$^3$ (soma de todos os níveis n $\geq$ 0, com a esfera inscrita em ABCD) & V = 2$\cdot$v0 = ($\pi$$\surd$6/108)$\cdot$a$^3$, com v0 = $\pi$$\surd$6$\cdot$a$^3$/216 & igual & ambos \\
\bottomrule\end{longtable}%
\begin{itemize}%
\item%
\textbf{Q10 --- volume total das esferas interiores ao tetraedro ABCD)} As duas dão a mesma resposta pelo mesmo caminho. H = a$\surd$6/3, r = H/4 = a$\surd$6/12. O plano paralelo a BCD e tangente à esfera fica à altura 2r = H/2 e passa pelos pontos médios das arestas AB, AC e AD. Cada novo tetraedro é semelhante ao anterior com razão 1/2. Por isso, a cada nível o número de esferas é multiplicado por 4 e o volume de cada esfera por 1/8. A soma é v0$\cdot$$\Sigma$(1/2)\^{}n = 2v0. Conferi em Python: v0 = 4$\pi$ r$^3$/3 = $\pi$$\surd$6a$^3$/216 $\approx$ 0,035626$\cdot$a$^3$; V = $\pi$$\surd$6a$^3$/108 $\approx$ 0,071253$\cdot$a$^3$; r$^3$ = $\surd$6a$^3$/288; V/V\_ABCD = $\pi$$\surd$3/9 $\approx$ 0,6046 < 1. A transcrição do enunciado bate com o PDF original (página 2). A resolução também prova que as esferas não se sobrepõem (Passo 5), o que o gabarito só supõe. Isso não muda o resultado.%
\end{itemize}%
\paragraph{Pontos de dúvida}%
\begin{itemize}%
\item%
\textit{(statement)} O enunciado pede o volume de 'todas as esferas interiores ao tetraedro ABCD'. A resolução e o gabarito incluem a esfera inscrita original (n = 0) e chegam a 2v0 = $\pi$$\surd$6a$^3$/108. Se alguém contasse só as esferas criadas pelo processo (n $\geq$ 1), o resultado seria v0 = $\pi$$\surd$6a$^3$/216. A redação favorece incluir a primeira esfera, e a resolução registra essa escolha de forma explícita.%
\item%
\textit{(key)} O gabarito soma os volumes sem justificar que as esferas não se sobrepõem. As esferas só se tocam em pontos: a esfera inscrita fica em $\lambda$\_X $\leq$ 1/2 e os filhos ficam em $\lambda$\_X $\geq$ 1/2. A conclusão está certa, mas falta esse passo, que a resolução demonstra no Passo 5.%
\end{itemize}%
\end{document}